% 原书PDF第21—25页，4.1.3.7全节；止于4.2标题前。
\subsubsection{利用二倍体基因组的 RNA-seq 数据估计等位基因特异性表达}
\bookdeepheading{paragraph}{4.1.3.7.1}{生成模型}

经常可以观察到两个等位基因的表达存在差异，即等位基因特异性表达（allele-specific expression，ASE）。RNA-seq 技术为识别 ASE 提供了有力工具。图4.10给出包含 ASE 的读段生成模型贝叶斯网络，本节大部分材料来自 Nariai et al.（2016）。模型包括四个变量：异构体选择变量 $G_n$、单倍型选择变量 $H_n$、起始位置变量 $S_n$，以及观测读段变量 $R_n$。$G_n=i$ 的定义同前，表示读段 $n$ 由第 $i$ 个异构体生成。

$H_n$ 表示生成读段 $n$ 的单倍型。设 $H_n=0$ 表示来自父源单倍型，$H_n=1$ 表示来自母源单倍型。$S_n$ 表示读段 $n$ 的起始位置。模型考虑两个参数向量：$\theta=[\theta_0,\theta_1,\ldots,\theta_M]^T$ 表示异构体丰度；$\phi=[\phi_0,\phi_1,\ldots,\phi_M]^T$ 表示每个异构体中父源单倍型的比例。假设
\[
\sum_{i=0}^{M}\theta_i=1,\qquad0\le\theta_i\le1,\qquad0\le\phi_i\le1.
\]
\begin{figure}[htbp]
\centering
\begin{tikzpicture}[x=1cm,y=1cm,>=Stealth,font=\small]
\node[ellipse,draw=inkblue,minimum width=3.2cm,minimum height=1.65cm] (ch04a10G) at (2,4.8) {};
\node at (2,4.3) {$G_n$};\node[text=red] at (.8,5.08) {$M$};
\draw[magenta] (1.2,5.4)--(2.1,5.4);\draw[green!65!black] (2.1,5.4)--(2.9,5.4);
\draw[magenta] (1.2,5.12)--(2.1,5.12);\draw[inkblue] (2.1,5.12)--(2.6,5.12);
\draw[magenta] (1.2,4.87)--(2.1,4.87);\draw[inkblue] (2.1,4.87)--(3.2,4.87);
\node at (-.5,4.8) {$\theta$};\draw[green!65!black,->] (-.2,4.8)--(ch04a10G.west);
\node[ellipse,draw=inkblue,minimum width=3.5cm,minimum height=1.85cm] (ch04a10H) at (2,.7) {};
\node at (2,1.3) {父源};\draw[red] (.9,1.05)--(3.1,1.05);
\node at (2,.6) {母源};\draw[green!65!black] (.9,.35)--(3.1,.35);
\node at (2,.05) {$H_n$};
\node at (-.5,.7) {$\phi$};\draw[green!65!black,->] (-.2,.7)--(ch04a10H.west);
\node[ellipse,draw=inkblue,minimum width=3.1cm,minimum height=1.75cm] (ch04a10S) at (5.5,2.65) {};
\node at (5.5,2.2) {$S_n$};
\draw[inkblue] (4.5,2.95)--(6.5,2.95);
\foreach \x in {4.5,4.95,5.4,5.85,6.3}{\fill[red] (\x,2.95) circle(.06);}
\draw[orange,->] (4.95,2.35)--(4.95,2.9);
\node[ellipse,draw=inkblue,minimum width=3.4cm,minimum height=1.9cm] (ch04a10R) at (10.3,2.5) {};
\node at (10.3,1.95) {$R_n$};
\draw[inkblue] (9.25,2.65)--(11.35,2.8);
\foreach \x in {9.3,9.75,10.2,10.65,11.1}{\fill[red] (\x,2.85) circle(.06);}
\node at (11,4.6) {读段序列};\draw[green!65!black,->,thin] (10.2,2.85)--(10.95,4.35);
\node at (5.2,6.2) {$N$ 条独立同分布的读段};
\draw[green!65!black,->] (ch04a10G)--(ch04a10H);
\draw[green!65!black,->] (ch04a10G)--(ch04a10S);
\draw[green!65!black,->] (ch04a10G)--(ch04a10R);
\draw[green!65!black,->] (ch04a10H)--(ch04a10S);
\draw[green!65!black,->] (ch04a10H)--(ch04a10R);
\draw[green!65!black,->] (ch04a10S)--(ch04a10R);
\end{tikzpicture}
\origfig{10}{ASE 的图形生成模型。}
\label{fig:ch04a-10}
\end{figure}


下面计算数据的完整似然。假设给定参数向量 $\theta$ 与 $\phi$ 后，图4.10各节点上的变量构成贝叶斯网络。根据贝叶斯网络的分解规则，数据似然可分解为各因子概率的乘积：
\begin{equation}
\begin{aligned}
L(\theta,\phi)&=P(G_n,H_n,S_n,R_n\mid\theta,\phi)\\
&=P(G_n\mid\theta)P(H_n\mid G_n,\phi)P(S_n\mid G_n,H_n)P(R_n\mid G_n,H_n,S_n).
\end{aligned}
\tag{4.42}
\end{equation}
为计算似然，先分别计算各组成部分。与前面相同，$P(G_n=i\mid\theta)$ 是读段 $n$ 由第 $i$ 个异构体生成的概率，因此定义 $P(G_n=i\mid\theta)=\theta_i$。

将 $P(H_n=0\mid G_n=i,\phi)$ 定义为 $\phi_i$，将 $P(H_n=1\mid G_n=i,\phi)$ 定义为 $1-\phi_i$。设 $l_{ih}$ 为单倍型 $h$ 上异构体 $i$ 的长度。概率 $P(S_n=j\mid G_n=i,H_n=h)$ 表示给定异构体 $i$ 和单倍型 $h$ 时，读段 $n$ 从位置 $j$ 生成的概率，计算为
\[
P(S_n=j\mid G_n=i,H_n=h)=\frac1{l_{ih}-L+1},
\]
其中 $L$ 为读段 $n$ 的长度。为便于讨论，定义指示变量
\begin{equation}
Z_{nihj}=\begin{cases}
1,&G_n=i,\ H_n=h,\ S_n=j,\\
0,&\text{其他情形}.
\end{cases}
\tag{4.43}
\end{equation}
令 $\pi_n$ 为读段 $n$ 全部可能比对的 $(i,h,j)$ 元组集合。式 (4.42) 中的概率可表示为
\begin{equation}
\begin{aligned}
P(R_n\mid G_n,H_n,S_n)&=P(R_n=r_n\mid Z_{nihj}=1)\\
&=\prod_{x=1}^{L}\operatorname{subst}(r_n[x],q_n[x],c_{ih}[x]),
\end{aligned}
\tag{4.44}
\end{equation}
其中 $\operatorname{subst}(r,q,c)$ 是纳入质量评分的替换矩阵；$x$ 为比对位置；$r_n[x]$ 为读段 $n$ 在位置 $x$ 的核苷酸；$q_n[x]$ 为该位置的质量评分；$c_{ih}[x]$ 为单倍型 $h$ 上异构体 $i$ 的 cDNA 参考序列的核苷酸。替换矩阵可利用 Phred 碱基质量评分估计，也可根据读段 $n$ 与参考 cDNA 序列的最佳比对来估计。显然，式 (4.42) 的模型似然 $L(\theta,\phi)$ 能够完整描述读段 $n$ 的生成。

\bookdeepheading{paragraph}{4.1.3.7.2}{ASE 估计的变分贝叶斯方法}

利用第4.1.3.6.3节介绍的变分贝叶斯方法，可以估计模型参数与 ASE。本节介绍 Nariai et al.（2016）提出的参数及 ASE 估计方法。

变分贝叶斯方法需要指定先验分布。仍以 Dirichlet 分布作为 $\theta$ 的先验：
\begin{equation}
P(\theta)=\frac{\Gamma\left(\sum_{i=0}^{M}\alpha_i\right)}
{\prod_{i=0}^{M}\Gamma(\alpha_i)}\theta_i^{\alpha_i-1},
\tag{4.45}
\end{equation}
其中 $\Gamma(\cdot)$ 为伽马函数。当 $\alpha_i-1\ge0$ 时，$\alpha_i-1$ 表示先验中分配给相应异构体的读段数；当 $\alpha_i-1<0$ 时，表示该异构体丰度等于零。容易证明
\begin{equation}
E[\theta_i]=\frac{\alpha_i}{\sum_{j=0}^{M}\alpha_j}.
\tag{4.46}
\end{equation}
参数 $\phi$ 的先验为 Beta 分布：
\begin{equation}
P(\phi_i)=\frac1{B(\beta_{i1},\beta_{i2})}\phi_i^{\beta_{i1}-1}(1-\phi_i)^{\beta_{i2}-1},
\tag{4.47}
\end{equation}
其中
\[
B(\beta_{i1},\beta_{i2})=\frac{\Gamma(\beta_{i1})\Gamma(\beta_{i2})}{\Gamma(\beta_{i1}+\beta_{i2})}
\]
是 Beta 函数；$\beta_{i1}>0$ 与 $\beta_{i2}>0$ 为 Beta 分布的两个形状参数。Beta 分布常用于群体遗传学中等位基因频率的分布建模，也是二项分布的共轭先验概率分布。参数 $\beta_{i1}$ 与 $\beta_{i2}$ 分别表示先验中分配给父源与母源单倍型的读段数。可以证明
\begin{equation}
E[\phi]=\frac{\beta_{i1}}{\beta_{i1}+\beta_{i2}}.
\tag{4.48}
\end{equation}
类似于结论4.3，可以构造下列估计 $\theta$ 和 $\phi$ 的变分贝叶斯算法（Nariai et al. 2016）。

\par\medskip\Needspace{5\baselineskip}
\noindent\textbf{结论4.5：ASE 估计的变分贝叶斯算法}\par
\noindent\textbf{步骤1：初始化。}
对每个转录本异构体，将 Dirichlet 分布参数初值设为 $\alpha_i^{(0)}=\alpha_0$，$i=1,\ldots,M$；将 Beta 分布参数初值设为 $\beta_{i1}^{(0)}=1$、$\beta_{i2}^{(0)}=1$，$i=1,\ldots,M$。

\noindent\textbf{步骤2：VBE 步。}
利用当前估计 $E_{\theta}[\theta^{(u)}]$ 和 $E_{\phi}[\phi^{(u)}]$，计算密度函数
\begin{equation}
q_Z^{(u+1)}(Z)=\prod_n\prod_i\prod_h\prod_k(\eta_{nihj})^{Z_{nihj}},
\tag{4.49}
\end{equation}
其中
\[
\eta_{n,i,h,j}=\frac{\rho_{n,i,h,j}}{\sum_{(i',h',j')\in\pi_n}\rho_{n',i',h',j'}}\qquad\text{或}
\]
\[
E_Z[Z_{nihj}]=\eta_{nihj},
\]
\[
\log\rho_{nihj}=
\begin{cases}
\begin{aligned}
&E_{\theta}[\log\theta_i^{(u)}]+E_{\phi}[\log\phi_i^{(u)}]+\log P(S_n\mid G_n)\\
&\qquad+\log P(R_n\mid G_n,H_n,S_n)
\end{aligned},&H_n=0,\\[1ex]
\begin{aligned}
&E_{\theta}[\log\theta_i^{(u)}]+E_{\phi}[\log(1-\phi_i^{(u)})]+\log P(S_n\mid G_n)\\
&\qquad+\log P(R_n\mid G_n,H_n,S_n)
\end{aligned},&H_n=1,
\end{cases}
\]
其中
\begin{equation}
E_{\theta}[\log\theta_i^{(u)}]
=\psi(\alpha_i^{(u)})-\psi\left(\sum_{j=1}^{M}\alpha_j^{(u)}\right),
\tag{4.50}
\end{equation}
\[
\psi(\alpha)=\frac{d\Gamma(\alpha)/d\alpha}{\Gamma(\alpha)}
\]
为双伽马函数，并且
\begin{equation}
E_{\phi}[\log\phi_i^{(u)}]
=\psi(\beta_{i1}^{(u)})-\psi(\beta_{i1}^{(u)}+\beta_{i2}^{(u)}),
\tag{4.51}
\end{equation}
以及
\[
E_{\phi}[\log(1-\phi_i^{(u)})]
=\psi(\beta_{i1}^{(u)})-\psi(\beta_{i1}^{(u)}+\beta_{i2}^{(u)}).
\]

\noindent\textbf{步骤3：VBM 步。}
利用当前估计 $q_Z^{(u+1)}(Z)$，计算
\begin{equation}
E_{\theta}[\theta_i^{(u+1)}]=\frac{\alpha_i^{(u+1)}}{\sum_{i=0}^{M}\alpha_i^{(u+1)}},
\tag{4.52}
\end{equation}
其中
\begin{equation}
\begin{aligned}
\alpha_i^{(u+1)}&=\alpha_i^{(u)}+\sum_{(n',i,h',j')\in\pi_n}E_Z[Z'_{n',i,h',j}]\\
&=\alpha_i^{(u)}+\sum_{(n',i,h',j')\in\pi_n}\eta_{n'ih'j'}.
\end{aligned}
\tag{4.53}
\end{equation}
同样，利用当前估计 $q^{(u+1)}(Z)$，计算
\begin{equation}
E_{\phi}[\phi_i]=\frac{\beta_{i1}^{(u+1)}}{\beta_{i1}^{(u+1)}+\beta_{i2}^{(u+1)}},
\tag{4.54}
\end{equation}
其中
\[
\beta_{i1}^{(u+1)}=\beta_{i1}^{(u)}+\sum_{n',t,h'=0,j'}E_Z[Z_{n',t,h',j'}],
\]
以及
\[
\beta_{i2}^{(u+1)}=\beta_{i2}^{(u)}+\sum_{n',t,h'=1,j'}E_Z[Z_{n',t,h',j'}].
\]
\noindent\textbf{步骤4：停止准则。}
若 $\lVert E_{\theta}[\theta^{(u+1)}-\theta^{(u)}]\rVert_2^2<\varepsilon$，则停止；否则返回步骤2。
\FloatBarrier
